% !PRINTED 198
\weberpage{198}
%<<CONT
gleich $2n-2$. Diese Werthe sollen die \gesperrt{Pole der Gruppe}
heissen. Die genaue Abzählung dieser Pole giebt uns die
wichtigsten Aufschlüsse über Zahl und Beschaffenheit der
möglichen endlichen Gruppen.

Wenn die Substitutionsco\"efficienten $\beta, \gamma$ gleich Null sind, so
ist $\Theta$ eine multiplicative Substitution
\[
\begin{pmatrix} \alpha, & 0 \\ 0, & \delta \end{pmatrix}.
\]

Damit diese Substitution von endlichem Grade sein kann,
muss der Quotient $\alpha : \delta = \varepsilon$ eine \gesperrt{Einheitswurzel sein,
deren Grad ein Theiler von $n$ ist}. Als Pole dieser Substitution
hat man $x = 0$ und $x = \infty$ anzusehen, die der Bedingung
\[
x = \varepsilon x
\]
genügen.

Nach (4) kann man aus $G$ eine isomorphe Gruppe
\[
L G L^{-1} = G'
\]
ableiten, wenn für $L$ eine beliebige lineare Substitution
\[
L = \begin{pmatrix} A, & B \\ C, & D \end{pmatrix}
\]
genommen wird. Die Pole dieser Gruppe erhält man, wenn
man auf die Pole von $G$ die Substitution $L$ anwendet. Wenn also
$a$ einer der Pole von $G$ ist, so ist
\[
a' = \frac{Aa + B}{Ca + D}
\]
der entsprechende Pol von $G'$.

Man kann die Substitutionsco\"efficienten $A, B, C, D$ so bestimmen,
dass drei der Pole von $G'$ vorgeschriebene Werthe erhalten.
Um z. B. den Polen $a, b, c$ von $G$ die Pole $0, \infty, 1$ von $G'$
entsprechen zu lassen, setze man
\[
L(x) = \frac{c-b}{c-a} \, \frac{x-a}{x-b}.
\]
Wenn $a$ und $b$ die Pole einer und derselben Substitution $\Theta$
von $G$ sind, so ist also die entsprechende Substitution von $G'$
multiplicativ.

Wir erhalten hieraus den Satz:

\par\noindent 2.\ \gesperrt{Zu jeder Gruppe linearer Substitutionen kann
man eine transformirte Gruppe finden, in der einer
beliebigen der gegebenen Substitutionen,}
%CONT>>
